% ===== PRÉAMBULE CANONIQUE — à coller VERBATIM en tête de CHAQUE .tex =====
\documentclass[11pt,a4paper]{article}
\usepackage[utf8]{inputenc}\usepackage[T1]{fontenc}\usepackage{lmodern}
\usepackage[french]{babel}
\usepackage{amsmath,amssymb,mathtools}\usepackage{icomma}
\usepackage{siunitx}\sisetup{locale=FR}  % \qty{}{}, \si{}, \ang{}
\usepackage{xcolor}\usepackage{colortbl}
\usepackage{tikz}\usepackage{pgfplots}\pgfplotsset{compat=1.18}
\usepackage[siunitx,europeanresistors]{circuitikz} % circuits (élec.)
\usepackage[most]{tcolorbox}\usepackage{enumitem}\usepackage{array,booktabs}
\usepackage[a4paper,margin=2.2cm]{geometry}
\usetikzlibrary{arrows.meta,calc,angles,quotes,positioning,patterns,%
  backgrounds,shapes.geometric,decorations.pathmorphing,decorations.markings,intersections}
\definecolor{primary}{RGB}{222,49,99}\definecolor{primarydark}{RGB}{150,22,55}
\definecolor{defcol}{RGB}{0,114,178}\definecolor{methcol}{RGB}{52,92,148}
\definecolor{excol}{RGB}{0,135,90}\definecolor{attncol}{RGB}{200,40,55}
\tcbset{boxbase/.style={enhanced,breakable,boxrule=0.9pt,arc=2.5pt,
  left=3mm,right=3mm,top=2.3mm,bottom=2.3mm,fonttitle=\bfseries,coltitle=white}}
% --- boîtes TUTORIEL (noms EXACTS, ne pas renommer) ---
\newtcolorbox{cle}[1][]{boxbase,colback=defcol!6,colframe=defcol,
  title={\textsc{À retenir}\ifx\relax#1\relax\relax\else\textmd{ : \itshape #1}\fi}}
\newtcolorbox{methode}[1][]{boxbase,colback=methcol!7,colframe=methcol,sharp corners,
  title={\textsc{Méthode}\ifx\relax#1\relax\relax\else\textmd{ --- \itshape #1}\fi}}
\newtcolorbox{exemplebox}[1][]{boxbase,colback=excol!5,colframe=excol,
  title={\textsc{Exemple}\ifx\relax#1\relax\relax\else\textmd{ : \itshape #1}\fi}}
\newtcolorbox{piege}[1][]{boxbase,colback=attncol!6,colframe=attncol,
  title={\textsc{Piège}\ifx\relax#1\relax\relax\else\textmd{ : \itshape #1}\fi}}
\newtcolorbox{remarque}[1][]{boxbase,colback=black!4,colframe=black!45,
  title={\textsc{Remarque}\ifx\relax#1\relax\relax\else\textmd{ : \itshape #1}\fi}}
% --- boîtes EXERCICE / CORRIGÉ (noms EXACTS, auto-comptées) ---
\newtcolorbox[auto counter]{exo}[1][]{boxbase,colback=primary!4,colframe=primary,
  title={\textsc{Exercice}~\thetcbcounter\ifx\relax#1\relax\relax\else\textmd{ --- \itshape #1}\fi}}
\newtcolorbox[auto counter]{corrige}[1][]{boxbase,colback=excol!4,colframe=excol,
  title={\textsc{Corrigé}~\thetcbcounter\ifx\relax#1\relax\relax\else\textmd{ --- \itshape #1}\fi}}
\newcommand{\vv}[1]{\vec{#1}}\newcommand{\ds}{\displaystyle}
\newcommand{\R}{\mathbb{R}}\newcommand{\N}{\mathbb{N}}\newcommand{\Z}{\mathbb{Z}}
% =========================================================================
\begin{document}
% THEME_TITLE: Les trois lois de Newton
\sisetup{per-mode=symbol}

\noindent La \textbf{dynamique} relie les \emph{forces} au \emph{mouvement} : elle explique
\emph{pourquoi} un corps accélère, ralentit ou reste immobile. Tout repose sur les
\textbf{trois lois de Newton}, valables dans un référentiel galiléen.

\begin{cle}[les trois lois de Newton]
\begin{enumerate}[label=\textbf{\arabic*.},leftmargin=1.7em,itemsep=3pt]
  \item \textbf{Principe d'inertie.} Si la résultante est nulle, l'accélération est nulle :
        \[ \sum\vv F_{\text{ext}}=\vv 0 \;\Longleftrightarrow\; \vv a=\vv 0
           \quad(\text{repos ou MRU}). \]
  \item \textbf{Principe fondamental (PFD).} Sinon, la résultante impose l'accélération :
        \[ \sum\vv F_{\text{ext}}=m\,\vv a \qquad\Longrightarrow\qquad a=\frac{F_{\text{rés}}}{m}. \]
        Le vecteur $\vv a$ est \emph{colinéaire} à la résultante : \textbf{même direction, même sens}.
  \item \textbf{Actions réciproques.} $\ \vv F_{A\to B}=-\,\vv F_{B\to A}$ : deux forces d'égale
        intensité, de même direction, de sens opposés, s'exerçant sur \emph{deux corps différents}.
\end{enumerate}
\end{cle}

\begin{methode}[résoudre un problème avec le PFD]
\begin{enumerate}[label=\arabic*),leftmargin=1.7em,topsep=2pt,itemsep=1.5pt]
  \item Choisir le \textbf{système} étudié et faire le \textbf{bilan des forces}.
  \item Calculer la \textbf{résultante} $\vv F_{\text{rés}}=\sum\vv F_{\text{ext}}$ (somme vectorielle).
  \item Si $\vv F_{\text{rés}}=\vv 0$ : $1^{\text{re}}$ loi, $a=0$ (repos ou MRU). Sinon : $a=F_{\text{rés}}/m$.
  \item Vérifier les unités : $\si{\newton}=\si{\kilogram\metre\per\second\squared}$, donc $a$ en \si{\metre\per\second\squared}.
\end{enumerate}
\end{methode}

\begin{center}
\begin{tikzpicture}[>=Stealth,scale=1]
  \draw[thick,black!55] (-3,0) -- (3.4,0);
  \fill[primarydark] (-1.4,0) rectangle (-0.4,0.9);
  \node[white] at (-0.9,0.45){$m$};
  \draw[->,line width=1.1pt,defcol] (-0.4,0.45) -- ++(1.9,0) node[right]{$\vv F_{\text{rés}}$};
  \draw[->,line width=1.1pt,primary] (-0.9,1.25) -- ++(1.4,0) node[right]{$\vv a$};
  \node[font=\small,align=center] at (-0.9,-0.55){$\vv F_{\text{rés}}$ et $\vv a$ : même direction, même sens};
\end{tikzpicture}
\end{center}

\begin{exemplebox}[même force, deux masses]
On pousse un chariot avec une force résultante de \qty{20}{\newton}.
Vide ($m=\qty{10}{\kilogram}$) : $a=\tfrac{20}{10}=\qty{2}{\metre\per\second\squared}$.
Plein ($m=\qty{20}{\kilogram}$) : $a=\tfrac{20}{20}=\qty{1}{\metre\per\second\squared}$.
Doubler la masse \emph{divise par deux} l'accélération : c'est l'inertie.
\end{exemplebox}

\begin{piege}[trois confusions classiques]
\begin{itemize}[leftmargin=1.4em,topsep=2pt,itemsep=1.5pt]
  \item « Forces équilibrées » ne veut pas dire « immobile » : cela veut dire \textbf{vitesse constante}
        (repos \emph{ou} MRU).
  \item Ce n'est pas l'accélération qui « crée » la force : la \textbf{force est la cause}, $\vv a$ l'effet.
  \item Les deux forces d'action--réaction \textbf{ne se compensent jamais} : elles agissent sur des
        corps \emph{différents}. On ne dit pas « la force d'un corps » mais « la force exercée \emph{par} A \emph{sur} B ».
\end{itemize}
\end{piege}

\begin{remarque}
Une force se mesure en \textbf{newtons} : $\qty{1}{\newton}$ communique $\qty{1}{\metre\per\second\squared}$ à une masse
de $\qty{1}{\kilogram}$. La masse $m$ (en \si{\kilogram}) mesure l'\emph{inertie} : plus elle est grande, plus le
corps résiste au changement de mouvement.
\end{remarque}
\end{document}
